% ==================================================
% CHAPTER 1: INTRODUCTION
% ==================================================

\chapter{Introduction}

\section{Background}

Students, developers and researchers read hundreds of web pages a week:
articles, documentation, forum threads, papers and videos. Much of what they
read is useful again later, but by then they usually remember only the idea,
not the title, the website or the exact words. Browser history records
\emph{that} a page was visited, and it can be searched only by title and
address. Bookmarks have to be made in advance and organised by hand. Web
search engines search the whole internet rather than the pages a person has
actually read.

MemoryLane treats the pages a person reads as a searchable personal memory.
It stores the readable text of each saved page, turns it into a
representation that supports search by meaning as well as by keywords, and
lets the user find a page again by describing it in their own words:

\begin{center}
\textit{``the article about reducing database query latency''} \quad
\textit{``the reddit post on rust lifetimes from last week''}
\end{center}

\section{Problem Statement}

Existing tools answer \emph{navigation} questions (``which pages did I
visit?'') but not \emph{content} questions (``where did I read about X?'').
A personal web memory has to solve several problems at once:

\begin{enumerate}
  \item \textbf{Capture} the useful text of a page, without menus, adverts
  and cookie banners, including pages such as videos whose content is not in
  the HTML.
  \item \textbf{Respect privacy}: never store an inbox, a bank statement or a
  private document, and keep every user's pages invisible to everyone else.
  \item \textbf{Search by meaning and by exact words}: a query may paraphrase
  the page, or may contain a precise name, error code or site.
  \item \textbf{Present pages, not fragments}: long pages must be split for
  search, but the user should see each page once, with the passage that
  matched.
  \item \textbf{Be fast and dependable} enough to use every day, on modest
  hardware.
\end{enumerate}

\section{Objectives}

The objectives set in the proposal, and refined after the mid-semester
evaluation, were:

\begin{enumerate}
  \item Build a browser extension that saves pages manually and, optionally,
  automatically, with privacy safeguards.
  \item Build a content pipeline: canonical URLs, normalisation, cleaning,
  chunking and embedding.
  \item Build hybrid retrieval combining BM25 and dense vectors, with
  page-level results and snippets.
  \item Make the system multi-user and secure, and deploy it publicly over
  HTTPS.
  \item Evaluate retrieval quality with standard information-retrieval
  metrics, and use the results to improve the system.
  \item Follow sound software-engineering practice: requirements, design
  records, automated tests, documentation and reproducible deployment.
\end{enumerate}

\section{From Prototype to Final Product}

The mid-semester report described a single-user prototype running locally.
The final product keeps its core idea and rebuilds almost every component.
Table~\ref{tab:prototype-vs-final} summarises the changes.

\begin{table}[!htbp]
\centering
\small
\begin{tabularx}{\textwidth}{>{\raggedright\arraybackslash}p{0.22\textwidth} L L}
\toprule
\textbf{Aspect} & \textbf{Prototype (mid-semester)} & \textbf{Final product} \\
\midrule
Users & One user, no sign-in & Multi-user, Amazon Cognito (Google or email), linked by verified email \\
Deployment & Local Docker & AWS EC2, Docker Compose, Caddy with HTTPS, \texttt{memorylane.site} \\
Embedding model & \texttt{all-MiniLM-L6-v2} & \texttt{bge-small-en-v1.5} with a query instruction \\
Chunking & 300--500 words & 150 words, 30-word overlap, title embedded with each chunk \\
Cleaning & Extension text & trafilatura main-content extraction, junk-line removal \\
Fusion & Min-max scores, 0.5\,/\,0.5 & Reciprocal Rank Fusion ($k=60$) \\
Query understanding & None & Site words, time phrases and URLs in the query \\
Result quality & Best chunk per page & Relevance floor, keyword-coverage check, sentence snippets \\
YouTube & Server fetched transcripts & Fetched in the user's browser (YouTube blocks servers) \\
Automatic capture & Saved most pages & Only pages actually read and that look like content; private pages refused twice \\
Library & None & Library page, delete with undo, delete all data \\
Support & None & In-app problem reports, emailed to the maintainer \\
Tests & A few scripts & 551 automated tests (415 backend, 136 frontend) \\
Documentation & README & 48-page documentation site, 11 decision records, meaning-based docs search \\
Evaluation & 3 metrics, 3 runs & 7 quality metrics, latency, empty-result rate, bootstrap intervals, 5 runs \\
\bottomrule
\end{tabularx}
\caption{The prototype and the final product}
\label{tab:prototype-vs-final}
\end{table}
